\section{数值数据的随机表示}\label{stochastic-representation-of-numerical-data}

随机计算机是一种增量计算机或"计数"计算机，其计算涉及无序逻辑电平序列的交互（而非二进制"字"之间的数字运算），这与数字微分分析仪\(^6\)、运算混合计算机\(^7\)以及相位计算机\(^8\)类似. 在所有这些计算机中，数值在存储时表示为二进制字，在计算时表示为时钟序列中 ON 逻辑电平所占比例（或运算计算机和异步 DDA 中的脉冲频率）. 常规计算机中，逻辑电平序列以确定性方式生成，通常呈现规律性模式或重复性. 而在随机计算机中，每个逻辑电平均由统计独立的过程生成，仅该过程的生成概率由待表示的数值决定.

图 1 展示了各类计算机中的典型序列，说明了这种区别：(a) 是同步速率乘法器的输出，或 DDA 中"R 寄存器"的输出，对应于存储量满 $\frac{3}{4}$ 的情况——注意 ON 和 OFF 逻辑电平分布均匀；(b) 是相位计算机中差分

\begin{figure}[htbp]
    \centering
    \begin{tabular}{ll}
        (a) & O I I O I I O I I O I I 速率乘法器输出 \\
        (b) & 差分存储器输出 \\
        (c) & O I O I I I I I O O I I I O \\
        (d) & \\ \\
        (e) & \\ \\
        (f) & 随机序列总体
    \end{tabular}
    \caption{各类增量计算机中的典型计算序列}
\end{figure}

存储器的输出——在一个周期内，所有的 OFF 逻辑电平都出现在 ON 逻辑电平之前，形成 mark/space 调制信号的同步等价形式；(c) 至 (f) 是生成概率为 $\frac{3}{4}$ 的随机序列示例——任何序列[包括 (a) 和 (b)]都可能出现，但在这类序列的大样本中，ON 电平比例服从均值为 $\frac{3}{4}$ 的二项分布.

尽管概率是能够表示模拟数据而无量化误差的连续变量，但该变量无法被精确测量，且会受到估计方差的影响. 比较各类计算机以 $1/N$ 精度承载模拟数据所需的电平数，可见方差对表示效率的影响：

\begin{itemize}
    \item 模拟计算机需要一个连续电平；
    \item 数字计算机需要 $\log_2 \mathbf{kN}$ 个有序二进制电平；
    \item DDA 需要 $\mathbf{kN}$ 个无序二进制电平；以及
    \item 随机计算机需要 $\mathbf{kN}^2$ 个无序二进制电平；
\end{itemize}

其中 $\mathbf{k} > 1$ 是表示舍入误差或方差影响的常数，例如 $\mathbf{k} = 10$. 随机计算机的 $\mathbf{N}^2$ 项源于：估计生成概率的期望误差随所采样序列长度的平方根而减小.

尽管从 $1 : \log_2\mathbf{N} : \mathbf{N} : \mathbf{N}^2$ 这一进展来看，随机计算机在数量表示方面效率最低，但随机序列缺乏相干性或规律性，使得可以使用简单的硬件来完成用这种形式表示的数据的复杂计算.
